% AUTO-GENERATED by scripts/make_tables.py -- DO NOT EDIT BY HAND
% Source: asi2_rank512_control_20260922T080042Z (results/asi2_rank512_control_20260922/)
\begin{table}[h]
\centering
\footnotesize
\caption{Rank-512 control (4 seeds, identical protocol
and run family as Table~\ref{tab:shuffled}, \texttt{asi2\_rank512\_control\_20260922T080042Z}).
At matched rank the causal core is marginally better than persistent
(paired causal$-$persistent $= -0.00132$, t(3)$=-65.94$;
the difference is statistically resolvable but is well inside any
practical non-inferiority margin). A persistent r512 factor beats the
causal r384 core by 14.6\% relative MSE
(paired t(3)$=-514$), so capacity dominates generation at
this scale.}
\label{tab:rank512control}
\begin{tabular}{lcc}
\toprule
Method (r512) & Rel.\ MSE & Cosine \\
\midrule
Static SVD & $0.7541\pm0.0017$ & $0.5003\pm0.0016$ \\
Persistent low rank & $0.3833\pm0.0012$ & $0.7857\pm0.0007$ \\
Shuffled diag core & $0.3820\pm0.0012$ & $0.7866\pm0.0007$ \\
\emph{Causal diag core} & \textbf{$0.3820\pm0.0012$} & \textbf{$0.7865\pm0.0007$} \\
\midrule
\multicolumn{3}{l}{Cross-rank (not parameter-matched at equal MiB):} \\
Persistent low rank r512 & $0.3833\pm0.0012$ & $0.7857\pm0.0007$ \\
Causal diag core r384 & $0.4489\pm0.0013$ & $0.7429\pm0.0009$ \\
\bottomrule
\end{tabular}
\end{table}
